\documentclass[14pt,a4paper]{extarticle}
% =============================================================
% preamble.tex — оформление отчётов по СТП БГУИР 01-2024
% =============================================================

% Шрифты: под pdflatex — Times через tempora + T2A,
% под xelatex/tectonic — системный Times New Roman (запасной вариант STIX Two Text)
\usepackage{iftex}
\ifPDFTeX
  \usepackage[utf8]{inputenc}
  \usepackage[T2A]{fontenc}
  \usepackage{tempora}
\else
  \usepackage{fontspec}
  \IfFontExistsTF{Times New Roman}
    {\setmainfont{Times New Roman}}
    {\setmainfont{STIX Two Text}}
  \IfFontExistsTF{DejaVu Sans Mono}
    {\setmonofont{DejaVu Sans Mono}[Scale=0.82]}
    {\setmonofont{DejaVuSansMono}[Extension=.ttf, BoldFont=*-Bold,
                                  ItalicFont=*-Oblique, Scale=0.82]}
\fi
\usepackage[russian]{babel}

% Поля: слева 30 мм, справа 15 мм, сверху и снизу 20 мм
\usepackage{geometry}
\geometry{left=30mm, right=15mm, top=20mm, bottom=20mm}

\usepackage{setspace}
\singlespacing

\pretolerance=1000
\tolerance=2000
\emergencystretch=3em

\usepackage{indentfirst}
\setlength{\parindent}{1.25cm}

% Номер страницы — снизу справа
\usepackage{fancyhdr}
\pagestyle{fancy}
\fancyhf{}
\fancyfoot[R]{\thepage}
\renewcommand{\headrulewidth}{0pt}
\fancypagestyle{plain}{%
  \fancyhf{}%
  \fancyfoot[R]{\thepage}%
  \renewcommand{\headrulewidth}{0pt}%
}

\usepackage{amsmath,amssymb}
\usepackage{graphicx}
\usepackage{float}
\usepackage{booktabs}
\usepackage{tabularx}
\usepackage{multirow}
\usepackage{longtable}
\usepackage{array}

\usepackage{caption}
\captionsetup[figure]{labelsep=endash, justification=centering,
                      font={normalsize}, name=Рисунок}
\captionsetup[table]{labelsep=endash, justification=raggedright,
                     singlelinecheck=false, font={normalsize}, name=Таблица}

\usepackage{titlesec}
\titleformat{\section}
  {\newpage\normalfont\fontsize{14}{18}\selectfont\bfseries}
  {\thesection}{0.5em}{\MakeUppercase}
\titlespacing*{\section}{1.25cm}{0pt}{14pt}
\titleformat{\subsection}
  {\normalfont\fontsize{14}{18}\selectfont\bfseries}
  {\thesubsection}{0.5em}{}
\titlespacing*{\subsection}{1.25cm}{14pt}{14pt}
\titleformat{\subsubsection}
  {\normalfont\fontsize{14}{18}\selectfont\bfseries}
  {\thesubsubsection}{0.5em}{}
\titlespacing*{\subsubsection}{1.25cm}{14pt}{14pt}

\newcommand{\sectioncenter}[1]{%
  \newpage
  \phantomsection
  \addcontentsline{toc}{section}{#1}
  \begin{center}\textbf{\MakeUppercase{#1}}\end{center}
  \vspace{14pt}
}

\usepackage{tocloft}
\renewcommand{\cftsecfont}{\normalfont}
\renewcommand{\cftsecpagefont}{\normalfont}
\renewcommand{\cftsubsecfont}{\normalfont}
\renewcommand{\cftsubsecpagefont}{\normalfont}
\renewcommand{\cftsecleader}{\cftdotfill{\cftdotsep}}
\renewcommand{\cftdotsep}{1}
% babel переопределяет заголовки после преамбулы, поэтому правим через \addto
\addto\captionsrussian{\renewcommand{\contentsname}{\hfill\textbf{СОДЕРЖАНИЕ}\hfill}}
\renewcommand{\cftaftertoctitle}{\hfill}

\usepackage{enumitem}
\setlist[itemize]{noitemsep, topsep=0pt, leftmargin=1.25cm, label=--}
\setlist[enumerate]{noitemsep, topsep=0pt, leftmargin=1.25cm}

% Листинги кода и вывода программы
\usepackage{listings}
\usepackage{fancyvrb}
\lstset{
  basicstyle=\ttfamily\footnotesize,
  frame=single,
  breaklines=true,
  showstringspaces=false,
  inputencoding=utf8,
  extendedchars=true,
  columns=fullflexible,
  keepspaces=true,
  captionpos=t
}
\ifPDFTeX
\lstset{
  literate={а}{{\cyra}}1 {б}{{\cyrb}}1 {в}{{\cyrv}}1 {г}{{\cyrg}}1
           {д}{{\cyrd}}1 {е}{{\cyre}}1 {ж}{{\cyrzh}}1 {з}{{\cyrz}}1
           {и}{{\cyri}}1 {й}{{\cyrishrt}}1 {к}{{\cyrk}}1 {л}{{\cyrl}}1
           {м}{{\cyrm}}1 {н}{{\cyrn}}1 {о}{{\cyro}}1 {п}{{\cyrp}}1
           {р}{{\cyrr}}1 {с}{{\cyrs}}1 {т}{{\cyrt}}1 {у}{{\cyru}}1
           {ф}{{\cyrf}}1 {х}{{\cyrh}}1 {ц}{{\cyrc}}1 {ч}{{\cyrch}}1
           {ш}{{\cyrsh}}1 {щ}{{\cyrshch}}1 {ъ}{{\cyrhrdsn}}1
           {ы}{{\cyrery}}1 {ь}{{\cyrsftsn}}1 {э}{{\cyrerev}}1
           {ю}{{\cyryu}}1 {я}{{\cyrya}}1 {ё}{{\cyryo}}1
           {А}{{\CYRA}}1 {Б}{{\CYRB}}1 {В}{{\CYRV}}1 {Г}{{\CYRG}}1
           {Д}{{\CYRD}}1 {Е}{{\CYRE}}1 {Ж}{{\CYRZH}}1 {З}{{\CYRZ}}1
           {И}{{\CYRI}}1 {Й}{{\CYRISHRT}}1 {К}{{\CYRK}}1 {Л}{{\CYRL}}1
           {М}{{\CYRM}}1 {Н}{{\CYRN}}1 {О}{{\CYRO}}1 {П}{{\CYRP}}1
           {Р}{{\CYRR}}1 {С}{{\CYRS}}1 {Т}{{\CYRT}}1 {У}{{\CYRU}}1
           {Ф}{{\CYRF}}1 {Х}{{\CYRH}}1 {Ц}{{\CYRC}}1 {Ч}{{\CYRCH}}1
           {Ш}{{\CYRSH}}1 {Щ}{{\CYRSHCH}}1 {Ъ}{{\CYRHRDSN}}1
           {Ы}{{\CYRERY}}1 {Ь}{{\CYRSFTSN}}1 {Э}{{\CYREREV}}1
           {Ю}{{\CYRYU}}1 {Я}{{\CYRYA}}1 {Ё}{{\CYRYO}}1
           {—}{{---}}1 {–}{{--}}1 {…}{{\ldots}}1
}
\fi

% Подписи листингов — в том же стиле, что таблицы: «Листинг X.Y — Название»
\captionsetup[lstlisting]{labelsep=endash, justification=raggedright,
                          singlelinecheck=false, font={normalsize}}
\renewcommand{\lstlistingname}{Листинг}
\AtBeginDocument{\numberwithin{lstlisting}{section}}

\lstdefinestyle{python}{language=Python, keywordstyle=\bfseries}
\lstdefinestyle{shell}{language=bash, keywordstyle=\mdseries}
\lstdefinestyle{plain}{language={}, basicstyle=\ttfamily\scriptsize}

% Абзац после листинга должен сохранять отступ 1,25 см (СТП 2.1.1)
\makeatletter
\AtBeginDocument{\let\@endpetrue\@endpefalse}
\makeatother

% Подпись к блоку verbatim: \captionof вне плавающего окружения
% глобально обнуляет \parindent, поэтому отступ восстанавливается вручную
\newcommand{\listingcaption}[2]{%
  \captionof{lstlisting}{#1}\label{#2}%
  \setlength{\parindent}{1.25cm}%
}

\usepackage{hyperref}
\hypersetup{colorlinks=true, linkcolor=black, urlcolor=blue, citecolor=black}

\numberwithin{figure}{section}
\numberwithin{table}{section}
\numberwithin{equation}{section}

% Титульный лист отчёта по лабораторной работе
\newcommand{\makelabtitle}[2]{%
  \begin{titlepage}
  \thispagestyle{empty}
  \begin{center}

  МИНИСТЕРСТВО ОБРАЗОВАНИЯ РЕСПУБЛИКИ БЕЛАРУСЬ

  \vspace{0.5em}
  Учреждение образования

  \vspace{0.5em}
  \textbf{<<Белорусский государственный университет информатики и радиоэлектроники>>}

  \vspace{2em}
  Факультет компьютерных систем и сетей

  \vspace{0.5em}
  Кафедра программного обеспечения информационных технологий
  \end{center}

  \vspace{3em}
  \begin{center}
  \textbf{ОТЧЕТ}

  \vspace{0.5em}
  по лабораторной работе №#1

  \vspace{0.5em}
  по дисциплине

  \vspace{0.5em}
  <<Информационные и технические средства защиты авторских прав>>

  \vspace{0.5em}
  <<#2>>
  \end{center}

  \vspace{5em}
  \begin{flushright}
  \begin{minipage}{0.5\textwidth}
  \textbf{Выполнил:}\\
  Лебедевич Артем Владимирович\\
  магистрант группы 556301

  \vspace{0.6em}
  \textbf{Преподаватель:}\\
  Ярмолик Вячеслав Николаевич\\
  доктор технических наук\\
  профессор кафедры ПОИТ
  \end{minipage}
  \end{flushright}

  \vfill
  \begin{center}Минск, 2026~г.\end{center}
  \end{titlepage}
}

% Титульный лист реферата
\newcommand{\makerefattitle}[1]{%
  \begin{titlepage}
  \thispagestyle{empty}
  \begin{center}

  МИНИСТЕРСТВО ОБРАЗОВАНИЯ РЕСПУБЛИКИ БЕЛАРУСЬ

  \vspace{0.5em}
  Учреждение образования

  \vspace{0.5em}
  \textbf{<<Белорусский государственный университет информатики и радиоэлектроники>>}

  \vspace{2em}
  Факультет компьютерных систем и сетей

  \vspace{0.5em}
  Кафедра программного обеспечения информационных технологий
  \end{center}

  \vspace{3em}
  \begin{center}
  \textbf{РЕФЕРАТ}

  \vspace{0.5em}
  по дисциплине

  \vspace{0.5em}
  <<Информационные и технические средства защиты авторских прав>>

  \vspace{1.5em}
  на тему

  \vspace{0.5em}
  \textbf{<<#1>>}
  \end{center}

  \vspace{5em}
  \begin{flushright}
  \begin{minipage}{0.5\textwidth}
  \textbf{Выполнил:}\\
  Лебедевич Артем Владимирович\\
  магистрант группы 556301

  \vspace{0.6em}
  \textbf{Преподаватель:}\\
  Ярмолик Вячеслав Николаевич\\
  доктор технических наук\\
  профессор кафедры ПОИТ
  \end{minipage}
  \end{flushright}

  \vfill
  \begin{center}Минск, 2026~г.\end{center}
  \end{titlepage}
}


\begin{document}

\makelabtitle{1}{Микросервис ввода и модификации данных о клиентах банка. Модуль <<Клиенты>>}

\tableofcontents

\section{Цель работы}

Разработать учебную подсистему ввода и модификации данных о клиентах условного банка:
спроектировать схему данных со справочниками, реализовать микросервис с REST-интерфейсом
на Java и Spring, web-клиент с формами <<Список клиентов>> и добавления, редактирования и
удаления клиента, обеспечить проверку вводимых значений дважды --- в браузере и в сервисе ---
и покрыть разработанное ПО модульными, интеграционными и функциональными тестами.
Дополнительно к заданию интерфейс и сообщения сервиса сделаны на трёх языках: русском,
английском и белорусском.

\section{Задание}

Работа выполнена по варианту №\,4. Номер варианта определяет состав полей клиента: вариант
чётный, поэтому в форму
входит поле <<Пол>> и не входят <<Место работы>> и <<Должность>>; $4 \bmod 3 = 1$, поэтому
из трёх полей <<Город прописки>>, <<Адрес прописки>> и <<Военнообязанный>> реализуется первое.
Итоговый состав полей приведён в таблице~\ref{tab:fields}.

\begin{table}[H]
\caption{Поля клиента для варианта №\,4}
\label{tab:fields}
\centering
\footnotesize
\begin{tabularx}{\textwidth}{lXcl}
\toprule
Поле & Тип & Обяз. & Контроль значения \\
\midrule
Фамилия, имя, отчество & текстовый & да & только буквы \\
Дата рождения & дата & да & существующая дата, не позднее сегодня \\
Пол & radio & да & мужской или женский \\
Серия паспорта & текстовый & да & две латинские буквы \\
№ паспорта & текстовый с маской & да & 7 цифр, уникален вместе с серией \\
Кем выдан & текстовый & да & не пусто \\
Дата выдачи & дата & да & не раньше даты рождения \\
Идентификационный номер & текстовый с маской & да & \texttt{3140301A001PB5}, уникален \\
Место рождения & текстовый & да & не пусто \\
Город проживания & список & да & значение справочника \\
Адрес проживания & текстовый & да & не пусто \\
Телефон домашний & текстовый с маской & нет & \texttt{293-88-44} \\
Телефон мобильный & текстовый с маской & нет & \texttt{+375 (29) 314-15-92} \\
E-mail & текстовый & нет & формат адреса \\
Город прописки & список & да & значение справочника \\
Семейное положение & список & да & значение справочника \\
Гражданство & список & да & значение справочника \\
Инвалидность & список & да & значение справочника \\
Пенсионер & checkbox & да & да или нет \\
Ежемесячный доход & денежный & нет & неотрицательная сумма \\
\bottomrule
\end{tabularx}
\end{table}

Требования задания:
\begin{enumerate}
\item разработать схему данных со справочниками для элементов типа <<Список>> и реализовать
её в СУБД, работающей в памяти; база заполняется автоматически при старте;
\item создать микросервис для REST-подключения к базе данных;
\item разработать web-клиент с формой <<Список клиентов>>, отсортированной по фамилиям,
и формой добавления, удаления и редактирования клиентов;
\item проверять значения каждого поля дважды --- на клиенте и в REST;
\item добавить в базу не менее пяти клиентов, включая себя;
\item протестировать ПО на WebDriver с базой H2 в памяти по пунктам программы защиты.
\end{enumerate}

\section{Проектирование}

\subsection{Архитектура}

Модуль реализован как самостоятельный микросервис: у сервиса \texttt{client-service} своя
база данных, собственный REST-интерфейс и собственный web-клиент, который раздаётся тем же
процессом как статический ресурс. Остальные модули банка (депозиты, кредиты) получают сведения
о клиентах только через REST, поэтому сервис можно развернуть, остановить и масштабировать
независимо от них. Условный банк в системе называется BankEt. Структура сервиса показана
на рисунке~\ref{fig:arch}.

\begin{figure}[H]
\centering
\begin{tikzpicture}[node distance=7mm and 12mm]
\node[box, minimum width=4.2cm, minimum height=1.25cm, fill=codekw!6] (ui)
  {Web-клиент (Vue 3)\\[-1pt]{\footnotesize формы, маски ввода,}\\[-3pt]{\footnotesize первичная валидация}};
\node[box, right=24mm of ui, minimum width=4.4cm] (ctl) {ClientController\\[-1pt]{\footnotesize REST, \texttt{@Valid}}};
\node[box, below=of ctl, minimum width=4.4cm] (svc) {ClientService\\[-1pt]{\footnotesize бизнес-правила, уникальность}};
\node[box, below=of svc, minimum width=4.4cm] (repo) {Spring Data JPA\\[-1pt]{\footnotesize репозитории}};
\node[db, below=7mm of repo] (h2) {H2 в памяти};
\node[box, right=14mm of svc, minimum width=3.4cm] (adv) {ApiExceptionHandler\\[-1pt]{\footnotesize ответы 400, 404, 409}};
\draw[flow, <->] (ui) -- node[above, note] {JSON / HTTP} (ctl);
\draw[flow] (ctl) -- (svc);
\draw[flow] (svc) -- (repo);
\draw[flow] (repo) -- (h2);
\draw[flow, dashed] (svc) -- node[above, note] {ошибки} (adv);
\draw[flow, dashed] (ctl.east) -| (adv.north);
\begin{scope}[on background layer]
\node[draw=black!40, dashed, rounded corners=3pt, fit=(ctl)(svc)(repo)(h2)(adv), inner sep=8pt,
      label={[note]above:{микросервис client-service, порт 8081}}] {};
\end{scope}
\end{tikzpicture}
\caption{Структура микросервиса <<Клиенты>>}
\label{fig:arch}
\end{figure}

Использованные технологии: Java~17, Spring Boot~3.5 (Spring MVC, Spring Data JPA, Bean
Validation), H2 в режиме in-memory, Vue~3 на стороне браузера, JUnit~5, Mockito, MockMvc и
Selenium WebDriver в тестах, Maven в качестве системы сборки. Библиотека Vue подключается из
webjar и раздаётся самим сервисом, поэтому для запуска не нужны ни Node.js, ни доступ в
интернет.

\subsection{Схема данных}

Каждое поле типа <<Список>> вынесено в отдельный справочник: \texttt{city} (используется
дважды --- для города проживания и города прописки), \texttt{marital\_status},
\texttt{citizenship} и \texttt{disability}. Таблица \texttt{client} ссылается на справочники
внешними ключами (рисунок~\ref{fig:er}). Пол хранится символом \texttt{M}/\texttt{F} с
ограничением \texttt{check}, поскольку набор значений фиксирован и справочник для него не нужен.
Наименование в каждом справочнике хранится на трёх языках: \texttt{name}~--- русское,
\texttt{name\_en} и \texttt{name\_be}~--- английское и белорусское.

\begin{figure}[H]
\centering
\begin{tikzpicture}
\node[entity, text width=6.4cm] (client) {\textbf{client}
  \nodepart{two}
  id \hfill PK\\ last\_name, first\_name, middle\_name\\ birth\_date, sex, birth\_place\\
  passport\_series, passport\_number \hfill UQ\\ issued\_by, issue\_date\\
  identification\_number \hfill UQ\\ residence\_city\_id \hfill FK\\ residence\_address\\
  home\_phone, mobile\_phone, email\\ registration\_city\_id \hfill FK\\ marital\_status\_id \hfill FK\\
  citizenship\_id \hfill FK\\ disability\_id \hfill FK\\ pensioner, monthly\_income};
\node[entity, text width=3.1cm, right=26mm of client.north east, anchor=north west] (city) {\textbf{city}\nodepart{two} id \hfill PK\\ name \hfill UQ\\ name\_en, name\_be};
\node[entity, text width=3.1cm, below=5mm of city] (ms) {\textbf{marital\_status}\nodepart{two} id \hfill PK\\ name \hfill UQ\\ name\_en, name\_be};
\node[entity, text width=3.1cm, below=5mm of ms] (cz) {\textbf{citizenship}\nodepart{two} id \hfill PK\\ name \hfill UQ\\ name\_en, name\_be};
\node[entity, text width=3.1cm, below=5mm of cz] (dis) {\textbf{disability}\nodepart{two} id \hfill PK\\ name \hfill UQ\\ name\_en, name\_be};
\draw[flow] (client.east |- city.west) -- node[above, note] {N : 1 (дважды)} (city.west);
\draw[flow] (client.east |- ms.west) -- node[above, note] {N : 1} (ms.west);
\draw[flow] (client.east |- cz.west) -- node[above, note] {N : 1} (cz.west);
\draw[flow] (client.east |- dis.west) -- node[above, note] {N : 1} (dis.west);
\end{tikzpicture}
\caption{Схема данных модуля <<Клиенты>>}
\label{fig:er}
\end{figure}

Схема создаётся сценарием \texttt{schema.sql}, справочники и шесть тестовых клиентов
(включая автора работы) загружаются сценарием \texttt{data.sql} при каждом старте приложения.
Hibernate работает в режиме \texttt{ddl-auto=validate}: он не создаёт таблицы сам, а сверяет
сущности с явно описанной схемой и не даст сервису стартовать при расхождении. Ограничения
уникальности вынесены в саму базу (листинг~\ref{lst:schema}) --- это последний рубеж защиты
от дублей, который срабатывает даже в обход сервиса.

\begin{lstlisting}[style=sql, caption={Фрагмент \texttt{schema.sql}: ссылки на справочники и уникальность}, label={lst:schema}]
create table client (
    id                    bigint generated by default as identity primary key,
    last_name             varchar(60)  not null,
    ...
    sex                   varchar(1)   not null check (sex in ('M', 'F')),
    residence_city_id     bigint       not null references city,
    registration_city_id  bigint       not null references city,
    marital_status_id     bigint       not null references marital_status,
    citizenship_id        bigint       not null references citizenship,
    disability_id         bigint       not null references disability,
    pensioner             boolean      not null,
    monthly_income        decimal(14, 2),
    constraint uq_client_passport unique (passport_series, passport_number),
    constraint uq_client_identification unique (identification_number)
);
\end{lstlisting}

\subsection{REST-интерфейс}

Сервис предоставляет шесть операций (таблица~\ref{tab:rest}). Данные передаются в JSON,
даты --- текстом в формате \texttt{ДД.ММ.ГГГГ}: это позволяет серверу самому решать, является
ли присланная строка датой, и возвращать ошибку конкретного поля, а не общий сбой разбора
запроса. Любая ошибка возвращается в едином виде --- общее сообщение и словарь
<<поле --- сообщение>>, по которому web-клиент подсвечивает поля формы. Язык сообщений и
наименований справочников определяется заголовком запроса \texttt{Accept-Language}.

\begin{table}[H]
\caption{REST-интерфейс микросервиса <<Клиенты>>}
\label{tab:rest}
\centering
\small
\begin{tabularx}{\textwidth}{llX}
\toprule
Метод & Адрес & Назначение и коды ответа \\
\midrule
GET & \texttt{/api/clients} & список клиентов, отсортированный по фамилии; 200 \\
GET & \texttt{/api/clients/\{id\}} & данные клиента; 200, 404 \\
POST & \texttt{/api/clients} & добавление; 201, 400 --- ошибка валидации, 409 --- дубль \\
PUT & \texttt{/api/clients/\{id\}} & изменение; 200, 400, 404, 409 \\
DELETE & \texttt{/api/clients/\{id\}} & удаление; 204, 404 \\
GET & \texttt{/api/dictionaries} & содержимое четырёх справочников для списков формы \\
\bottomrule
\end{tabularx}
\end{table}

\section{Реализация}

\subsection{Структура проекта}

Сервис является модулем многомодульного Maven-проекта банка. Код разложен по слоям
(листинг~\ref{lst:tree}); общие для всех web-клиентов стили, маски ввода, выбор языка и
обёртка над \texttt{fetch} вынесены в модуль \texttt{ui-kit}, а серверная часть локализации
--- в пакет \texttt{i18n} общей библиотеки \texttt{bank-common}.

\listingcaption{Структура модуля \texttt{client-service}}{lst:tree}
\begin{Verbatim}[fontsize=\footnotesize, frame=single, rulecolor=\color{codeframe}]
src/main/java/by/bsuir/bank/client/
  domain/        Client, справочники City, MaritalStatus, Citizenship, Disability
  repository/    репозитории Spring Data JPA
  dto/           ClientRequest (форма с правилами валидации), ClientResponse, ApiError
  validation/    Masks — маски полей, Dates и @DateText — строгая проверка дат
  service/       ClientService — уникальность, справочники, сохранение
  web/           ClientController, DictionaryController, ApiExceptionHandler
src/main/resources/
  schema.sql     схема данных          data.sql   справочники и тестовые клиенты
  messages.properties, messages_en.properties, messages_be.properties — тексты сервиса
  static/        index.html, app.js — web-клиент на Vue 3, i18n.js — его словарь
src/test/java/   модульные (*Test), интеграционные и функциональные (*IT) тесты
\end{Verbatim}

\subsection{Web-клиент}

Web-клиент --- одностраничное приложение на Vue~3 с двумя формами. Форма <<Список
клиентов>> запрашивает \texttt{GET /api/clients} и показывает таблицу; сортировку по фамилии
выполняет сервис, чтобы порядок был одинаковым для любого потребителя REST. Форма клиента
используется и для добавления, и для редактирования; удаление доступно из списка и из формы.

Форма построена по описанию полей: для каждого поля в одном месте заданы тип,
обязательность, маска ввода и функция проверки, а разметка и валидация генерируются по этому
описанию; подписи берутся из словаря по имени поля. Добавление поля сводится к одной строке, и расхождение между тем, что нарисовано,
и тем, что проверяется, становится невозможным. Маска (\texttt{\#}~--- цифра, \texttt{A}~---
латинская буква) применяется при каждом нажатии клавиши: недопустимые символы отбрасываются,
разделители расставляются автоматически, поэтому в поле даты физически нельзя набрать текст.

\subsection{Двойная проверка маски и обязательности}

Проверка значений выполняется дважды. Первичная проверка в браузере нужна для удобства:
пользователь сразу видит сообщения у полей и не ждёт ответа сервера. Но она не защищает
базу --- запрос можно отправить в обход формы, поэтому сервис повторяет все проверки, и
именно он является источником истины.

Первичная проверка на клиенте показана в листинге~\ref{lst:client-validation}. Полужирным
выделен фрагмент, который для каждого поля проверяет обязательность (признак
\texttt{required}) и соответствие маске (функция \texttt{check} с регулярным выражением).

\begin{lstlisting}[style=js, style=marked, caption={Первичная проверка в web-клиенте (\texttt{static/app.js})}, label={lst:client-validation}]
const SECTIONS = [
  ...
  { title: 'passport', fields: [
    { name: 'passportSeries', @b@required: true, mask: 'AA'@/b@, placeholder: 'MP',
      @b@check: v => /^[A-Z]{2}$/.test(v) || t('check.passportSeries')@/b@ },
    { name: 'passportNumber', @b@required: true, mask: '#######'@/b@, placeholder: '1234567',
      @b@check: v => /^\d{7}$/.test(v) || t('check.passportNumber')@/b@ },
    { name: 'identificationNumber', @b@required: true, mask: '#######A###AA#'@/b@, placeholder: '3140301A001PB5',
      @b@check: v => /^\d{7}[A-Z]\d{3}[A-Z]{2}\d$/.test(v) || t('check.identificationNumber')@/b@ },
    ...
];

/* Проверка формы перед отправкой: возвращает { имяПоля: сообщение }. Пустой объект — ошибок нет. */
function validate(form) {
  const errors = {};
  @b@for (const field of FIELDS) {
    const raw = form[field.name];
    const value = typeof raw === 'string' ? raw.trim() : raw;
    const empty = value === null || value === undefined || value === '';
    if (empty) {
      if (field.required) errors[field.name] = t('common.required');
      continue;
    }
    const result = field.check ? field.check(value, form) : true;
    if (result !== true) errors[field.name] = result;
  }@/b@
  return errors;
}
\end{lstlisting}

Проверка в сервисе (листинг~\ref{lst:server-validation}) описана декларативно аннотациями
Bean Validation на записи \texttt{ClientRequest}; контроллер помечает параметр аннотацией
\texttt{@Valid}, и до бизнес-логики запрос с ошибками не доходит. Полужирным выделены
ограничения обязательности (\texttt{@NotBlank}, \texttt{@NotNull}) и маски (\texttt{@Pattern},
\texttt{@DateText}). Регулярные выражения масок собраны в классе \texttt{Masks} и по
содержанию совпадают с выражениями web-клиента. Сообщения заданы ключами в фигурных скобках:
текст подставляется из словаря на языке запроса.

\begin{lstlisting}[style=java, style=marked, caption={Серверная проверка: запись \texttt{ClientRequest} и маски \texttt{Masks}}, label={lst:server-validation}]
public record ClientRequest(

        @b@@NotBlank(message = REQUIRED)@/b@ @Size(max = 60, message = TOO_LONG)
        @b@@Pattern(regexp = Masks.NAME, message = LETTERS_ONLY)@/b@
        String lastName,
        ...
        @b@@NotBlank(message = REQUIRED) @DateText@/b@
        String birthDate,

        @b@@NotBlank(message = REQUIRED)
        @Pattern(regexp = Masks.PASSPORT_SERIES, message = "{validation.passportSeries}")@/b@
        String passportSeries,

        @b@@NotBlank(message = REQUIRED)
        @Pattern(regexp = Masks.PASSPORT_NUMBER, message = "{validation.passportNumber}")@/b@
        String passportNumber,
        ...
        @b@@NotBlank(message = REQUIRED)
        @Pattern(regexp = Masks.IDENTIFICATION_NUMBER, message = "{validation.identificationNumber}")@/b@
        String identificationNumber,
        ...
        @b@@Pattern(regexp = Masks.MOBILE_PHONE, message = "{validation.mobilePhone}")@/b@
        String mobilePhone,

        @b@@NotNull(message = REQUIRED)@/b@
        Long residenceCityId,
        ...
) { ... }

public final class Masks {
    public static final String NAME = "^[А-Яа-яЁёA-Za-z]+(?:[-' ][А-Яа-яЁёA-Za-z]+)*$";
    @b@public static final String PASSPORT_SERIES = "^[A-Z]{2}$";
    public static final String PASSPORT_NUMBER = "^\\d{7}$";
    public static final String IDENTIFICATION_NUMBER = "^\\d{7}[A-Z]\\d{3}[A-Z]{2}\\d$";
    public static final String HOME_PHONE = "^$|^\\d{3}-\\d{2}-\\d{2}$";
    public static final String MOBILE_PHONE = "^$|^\\+375 \\((25|29|33|44)\\) \\d{3}-\\d{2}-\\d{2}$";@/b@
}

@PostMapping
@ResponseStatus(HttpStatus.CREATED)
public ClientResponse create(@b@@Valid@/b@ @RequestBody ClientRequest request) {
    return service.create(request);
}
\end{lstlisting}

Три детали серверной проверки заслуживают отдельного пояснения.

\textbf{Пробел вместо имени.} Компактный конструктор записи отбрасывает пробелы по краям
всех строк ещё до проверок, поэтому строка из одного пробела превращается в пустую и
отклоняется как незаполненное обязательное поле.

\textbf{Несуществующая дата.} Стандартный разбор даты в Java по умолчанию работает в режиме
\texttt{SMART} и молча превращает 31~февраля в 28-е. Для строгой проверки формат создаётся с
\texttt{ResolverStyle.STRICT} (листинг~\ref{lst:dates}): и 31.02.2015, и 29.02.2015
отвергаются, а 29.02.2016 принимается. Собственное ограничение \texttt{@DateText}
дополнительно требует, чтобы дата лежала между 01.01.1900 и сегодняшним днём.

\begin{lstlisting}[style=java, caption={Строгий разбор даты (\texttt{validation/Dates.java})}, label={lst:dates}]
public static final DateTimeFormatter FORMAT =
        DateTimeFormatter.ofPattern("dd.MM.uuuu").withResolverStyle(ResolverStyle.STRICT);

public static Optional<LocalDate> parse(String text) {
    if (text == null) {
        return Optional.empty();
    }
    try {
        return Optional.of(LocalDate.parse(text.trim(), FORMAT));
    } catch (DateTimeParseException e) {
        return Optional.empty();
    }
}
\end{lstlisting}

\textbf{Необязательные поля.} Для телефонов и e-mail маска допускает пустую строку
(\texttt{\textasciicircum\$|...}), а сервис сохраняет пустое значение как \texttt{NULL}:
запись клиента только с обязательными полями выполняется без ошибок.

\subsection{Уникальность паспорта и идентификационного номера}

Маска не может обнаружить повторный ввод того же клиента --- это правило проверяется в
сервисе обращением к базе (листинг~\ref{lst:unique}). При редактировании собственная запись
клиента из сравнения исключается, иначе его нельзя было бы сохранить без смены паспорта.
Нарушение возвращается с кодом 409 и сообщениями у полей паспорта и идентификационного
номера. Если два одинаковых запроса придут одновременно и оба пройдут проверку, второй
будет остановлен ограничением уникальности в базе; обработчик ошибок превращает
\texttt{DataIntegrityViolationException} в тот же ответ 409.

\begin{lstlisting}[style=java, caption={Проверка уникальности (\texttt{service/ClientService.java})}, label={lst:unique}]
private void checkUnique(ClientRequest request, Long selfId) {
    Map<String, String> errors = new LinkedHashMap<>();
    clients.findByPassportSeriesAndPassportNumber(request.passportSeries(), request.passportNumber())
            .filter(other -> !other.getId().equals(selfId))
            .ifPresent(other -> {
                errors.put("passportSeries", Messages.get("client.duplicatePassport"));
                errors.put("passportNumber", Messages.get("client.duplicatePassport"));
            });
    clients.findByIdentificationNumber(request.identificationNumber())
            .filter(other -> !other.getId().equals(selfId))
            .ifPresent(other -> errors.put("identificationNumber", Messages.get("client.duplicateIdentification")));
    if (!errors.isEmpty()) {
        throw new FieldErrorException(HttpStatus.CONFLICT, Messages.get("client.duplicate"), errors);
    }
}
\end{lstlisting}

\subsection{Локализация}

Язык выбирается переключателем в шапке любого web-клиента. Выбор хранится в cookie
\texttt{bank\_lang}: cookie не зависит от порта, поэтому язык един для всех web-клиентов
системы. Перевод выполняется в трёх местах (таблица~\ref{tab:i18n}).

\begin{table}[H]
\caption{Где переводятся тексты}
\label{tab:i18n}
\centering
\small
\begin{tabularx}{\textwidth}{lXX}
\toprule
Что & Где хранится перевод & Как выбирается язык \\
\midrule
Подписи, кнопки, первичная проверка & словарь \texttt{i18n.js} web-клиента & функция \texttt{Bank.t} по cookie \\
Сообщения сервиса & \texttt{messages*.properties} & заголовок \texttt{Accept-Language} \\
Значения справочников & столбцы \texttt{name}, \texttt{name\_en}, \texttt{name\_be} & заголовок \texttt{Accept-Language} \\
\bottomrule
\end{tabularx}
\end{table}

Web-клиент при каждом REST-вызове передаёт выбранный язык в заголовке
\texttt{Accept-Language} (листинг~\ref{lst:i18n-js}). В сервисе язык запроса определяет
стандартный \texttt{AcceptHeaderLocaleResolver} с тремя поддерживаемыми языками и русским по
умолчанию, а тексты выдаёт \texttt{MessageSource} из файлов \texttt{messages.properties}
(русский), \texttt{messages\_en.properties} и \texttt{messages\_be.properties}. Тот же
источник использует Bean Validation, поэтому ключ \texttt{\{validation.required\}} в
аннотации превращается в <<Обязательное поле>>, <<Required field>> или <<Абавязковае поле>>.

\begin{lstlisting}[style=js, caption={Выбор текста и передача языка сервису (\texttt{ui-kit/bank.js})}, label={lst:i18n-js}]
lang: (document.cookie.match(/(?:^|; )bank_lang=(ru|en|be)/) || [])[1] || 'ru',

t(key, ...args) {
  const text = Bank.messages[Bank.lang][key] ?? Bank.messages.ru[key] ?? key;
  return text.replace(/\{(\d+)\}/g, (match, index) => args[index] ?? '');
},

async api(method, url, body) {
  const options = { method, headers: { 'Accept': 'application/json', 'Accept-Language': Bank.lang } };
  ...
\end{lstlisting}

Значения справочников --- это данные, а не тексты программы, поэтому их переводы хранятся
в базе рядом с русским наименованием. Сущность справочника отдаёт наименование на языке
запроса (листинг~\ref{lst:i18n-java}), и в ответе REST всегда одно поле \texttt{name}:
потребителю не нужно знать о существовании переводов. Данные самого клиента --- ФИО, адрес ---
не переводятся.

\begin{lstlisting}[style=java, caption={Наименование справочника на языке запроса}, label={lst:i18n-java}]
// DictionaryItem — общий предок справочников
public String getName() {
    return Messages.pick(name, nameEn, nameBe);
}

// Messages — тексты на языке пользователя
public static String pick(String ru, String en, String be) {
    String language = locale().getLanguage();
    String value = "en".equals(language) ? en : "be".equals(language) ? be : ru;
    return value == null || value.isBlank() ? ru : value;
}
\end{lstlisting}

\section{Тестирование}

\subsection{Уровни тестов}

Тесты образуют пирамиду из трёх уровней (таблица~\ref{tab:tests}). Модульные тесты
запускаются на фазе \texttt{test} и не поднимают Spring; интеграционные и функциональные
(имена классов оканчиваются на \texttt{IT}) выполняются плагином Failsafe на фазе
\texttt{verify}. Все уровни используют базу H2 в памяти; у каждого класса тестов своя база,
поэтому тесты не влияют друг на друга.

\begin{table}[H]
\caption{Тесты микросервиса <<Клиенты>>}
\label{tab:tests}
\centering
\small
\begin{tabularx}{\textwidth}{lcX}
\toprule
Класс тестов & Тестов & Что проверяется \\
\midrule
\multicolumn{3}{l}{\textit{Модульные --- без Spring и базы данных}} \\
\texttt{DateTextValidatorTest} & 16 & существующие и несуществующие даты \\
\texttt{ClientRequestValidationTest} & 49 & маски и обязательность каждого поля \\
\texttt{ClientServiceTest} & 8 & уникальность и справочники на заглушках Mockito \\
\midrule
\multicolumn{3}{l}{\textit{Интеграционные --- сервис целиком и H2}} \\
\texttt{ClientApiIT} & 39 & REST через MockMvc, серверная валидация, состояние таблицы \\
\texttt{ClientLanguageIT} & 9 & сообщения и справочники на трёх языках \\
\midrule
\multicolumn{3}{l}{\textit{Функциональные --- браузер через WebDriver}} \\
\texttt{ClientUiIT} & 37 & сценарии программы защиты и переключение языка в web-клиенте \\
\midrule
Всего & 158 & \\
\bottomrule
\end{tabularx}
\end{table}

\subsection{Программа защиты в виде автоматических тестов}

Функциональные тесты повторяют пункты, по которым проводится защита работы
(таблица~\ref{tab:defense}). Под данными <<Иванов Иван Иванович>> в задании понимаются данные
студента; в тестах пунктов 2 и 3 новый клиент вводится с паспортом и идентификационным
номером клиента Лебедевич из начального заполнения базы. Каждый тест после действий в
браузере выполняет SQL-запрос к таблице \texttt{client} --- проверку <<прямо в базе>>.

\begin{table}[H]
\caption{Соответствие пунктов защиты и тестов}
\label{tab:defense}
\centering
\small
\begin{tabularx}{\textwidth}{cXl}
\toprule
Пункт & Сценарий и ожидаемый результат & Тест \texttt{ClientUiIT} \\
\midrule
1 & те же данные ещё раз --- ошибка, в базе один клиент & \texttt{sameClientTwice} \\
2 & другой клиент с тем же паспортом --- ошибка & \texttt{anotherClientWithSamePassport} \\
3 & другой клиент с тем же идентификационным номером --- ошибка & \texttt{...SameIdentificationNumber} \\
4 & фамилия <<1234>> --- ошибка, фамилии 1234 в базе нет & \texttt{digitsInsteadOfLastName} \\
5 & пробел вместо имени --- ошибка, клиент не записан & \texttt{spaceInsteadOfFirstName} \\
6 & не дата, 31.02.2015, 29.02.2015; все поля с маской & \texttt{birthDateIsNotDate}, \texttt{maskedFields} \\
7 & нет одного обязательного поля (17 вариантов) --- ошибка & \texttt{requiredFieldIsMissing} \\
8 & только обязательные поля --- клиент записан & \texttt{onlyRequiredFields} \\
9 & двойная проверка на клиенте и на сервере & листинги~\ref{lst:client-validation}, \ref{lst:server-validation} \\
\bottomrule
\end{tabularx}
\end{table}

Поскольку web-клиент не пропускает некорректные данные до сервера, серверная валидация
проверяется отдельно: интеграционные тесты \texttt{ClientApiIT} отправляют те же неверные
данные прямо в REST, минуя браузер, и убеждаются, что сервис отвечает кодом 400 или 409 и
ничего не пишет в базу. Пример теста каждого уровня приведён в листингах~\ref{lst:it}
и~\ref{lst:ui}.

\begin{lstlisting}[style=java, caption={Интеграционный тест пункта 6 (\texttt{ClientApiIT})}, label={lst:it}]
@ParameterizedTest(name = "дата рождения «{0}»")
@ValueSource(strings = {"не дата", "31.02.2015", "29.02.2015"})
@DisplayName("п.6 Текст, не являющийся датой, — ошибка")
void birthDateIsNotDate(String birthDate) throws Exception {
    create(TestClients.with("birthDate", birthDate))
            .andExpect(status().isBadRequest())
            .andExpect(jsonPath("$.fields.birthDate").exists());

    assertThat(count("last_name = 'Иванов'")).isZero();
}

/** Проверка «прямо в базе»: SQL-запрос к таблице client в обход сервиса и репозиториев. */
private int count(String where) {
    return jdbc.queryForObject("select count(*) from client where " + where, Integer.class);
}
\end{lstlisting}

\begin{lstlisting}[style=java, caption={Функциональный тест пункта 2 через WebDriver (\texttt{ClientUiIT})}, label={lst:ui}]
@Test
@DisplayName("п.2 Другой клиент с паспортом студента — сообщение об ошибке")
void anotherClientWithSamePassport() {
    Map<String, Object> other = TestClients.valid();
    other.put("passportSeries", "MP");
    other.put("passportNumber", "3141592");          // паспорт клиента Лебедевич из начальных данных
    fill(other);
    save();

    assertThat(error("passportSeries")).isEqualTo("Клиент с таким паспортом уже зарегистрирован");
    assertThat(count("passport_series = 'MP' and passport_number = '3141592'")).isEqualTo(1);
    assertThat(count("last_name = 'Иванов'")).isZero();
}
\end{lstlisting}

\subsection{Результаты прогона}

Команда \texttt{./mvnw -pl client-service -am verify} собирает сервис и выполняет все
158~тестов; функциональные тесты запускают Chrome без окна, драйвер для него загружает
Selenium Manager. Итог прогона приведён в листинге~\ref{lst:run}.

\listingcaption{Результат выполнения тестов}{lst:run}
\VerbatimInput[fontsize=\fontsize{7}{8.8}\selectfont, frame=single, rulecolor=\color{codeframe}, samepage=true]{listings/tests.txt}

\section{Результаты работы}

После запуска сервис доступен по адресу \texttt{http://localhost:8081}. Форма <<Список
клиентов>> с начальным заполнением базы показана на рисунке~\ref{fig:list}: шесть клиентов
отсортированы по фамилии. В правой части шапки находится переключатель языка.

\screenshot{img/list.png}{Форма <<Список клиентов>>}{fig:list}

Форма добавления клиента, заполненная корректными данными, приведена на
рисунке~\ref{fig:form}. Обязательные поля отмечены звёздочкой, списки заполнены из
справочников, значения полей с маской набраны с автоматической расстановкой разделителей.

\screenshot[0.92\textwidth]{img/form.png}{Форма добавления клиента}{fig:form}

На рисунке~\ref{fig:errors} показана реакция первичной проверки на ошибки из программы
защиты: цифры в фамилии, пробел вместо имени, дата 31.02.2015, неполный номер паспорта,
мобильный телефон с недопустимым кодом и неполный e-mail. Запрос на сервер при этом не
отправляется.

\screenshot[0.92\textwidth]{img/form_errors.png}{Сообщения первичной проверки в браузере}{fig:errors}

Ошибку повторного ввода обнаруживает только сервер. На рисунке~\ref{fig:duplicate} новый
клиент введён с паспортом и идентификационным номером, которые уже принадлежат клиенту из
базы: сервис ответил кодом 409, а web-клиент показал его сообщения у соответствующих полей.

\screenshot[0.92\textwidth]{img/form_duplicate.png}{Сообщения сервиса о нарушении уникальности}{fig:duplicate}

После исправления данных клиент сохраняется и появляется в списке на своём месте по
алфавиту (рисунок~\ref{fig:added}).

\screenshot{img/list_added.png}{Список после добавления клиента}{fig:added}

Состояние таблицы можно проверить напрямую в консоли H2 по адресу \texttt{/h2-console}
(JDBC URL \texttt{jdbc:h2:mem:clients}). На рисунке~\ref{fig:h2} видно, что после всех
попыток в таблице семь записей и нет ни одинаковых паспортов, ни одинаковых идентификационных
номеров, ни клиента с фамилией <<1234>>.

\screenshot{img/h2.png}{Таблица \texttt{client} в консоли H2}{fig:h2}

Тот же интерфейс на других языках показан на рисунках~\ref{fig:en} и~\ref{fig:be}. На
английском выведен список клиентов: переведены заголовки, кнопки и название города из
справочника, а ФИО остаются данными клиента. На белорусском показана форма с ошибками:
сообщения первичной проверки взяты из словаря web-клиента.

\screenshot{img/list_en.png}{Список клиентов на английском языке}{fig:en}

\screenshot[0.92\textwidth]{img/form_be.png}{Форма клиента на белорусском языке}{fig:be}

\sectioncenter{Заключение}

В ходе работы разработан микросервис <<Клиенты>> на Java и Spring Boot с базой данных H2
в памяти, REST-интерфейсом и web-клиентом на Vue~3. Схема данных содержит таблицу клиентов и
четыре справочника для полей типа <<Список>>; база создаётся и заполняется автоматически при
старте приложения.

Корректность данных обеспечивается на трёх уровнях. Web-клиент применяет маски ввода и
проверяет форму до отправки, сервис повторяет проверку масок и обязательности средствами
Bean Validation и контролирует уникальность паспорта и идентификационного номера, а база
данных хранит ограничения уникальности и ссылочной целостности. Такое разделение
соответствует практике реальных банковских систем: проверка в браузере служит удобству
пользователя, а гарантии целостности даёт только сервер.

Интерфейс и сообщения сервиса доступны на русском, английском и белорусском языках:
тексты web-клиента переводятся в браузере, сообщения сервиса и значения справочников --- на
сервере по заголовку \texttt{Accept-Language}.

Работоспособность подтверждена 158 автоматическими тестами: 73 модульными, 48
интеграционными и 37 функциональными, которые через WebDriver проходят все пункты программы
защиты и сверяют результат с содержимым таблицы в базе.

\sectioncenter{Список использованных источников}

\begin{enumerate}
\item Уоллс, К. Spring в действии / К.~Уоллс. --- 6-е изд. --- М.: ДМК Пресс, 2022. --- 544~с.
\item Spring Boot Reference Documentation [Электронный ресурс]. --- Режим доступа:
\url{https://docs.spring.io/spring-boot/}. --- Дата доступа: 01.10.2026.
\item Jakarta Bean Validation 3.0 Specification [Электронный ресурс]. --- Режим доступа:
\url{https://jakarta.ee/specifications/bean-validation/3.0/}. --- Дата доступа: 01.10.2026.
\item Selenium WebDriver Documentation [Электронный ресурс]. --- Режим доступа:
\url{https://www.selenium.dev/documentation/webdriver/}. --- Дата доступа: 01.10.2026.
\item Vue.js Guide [Электронный ресурс]. --- Режим доступа: \url{https://vuejs.org/guide/}. ---
Дата доступа: 01.10.2026.
\end{enumerate}

\appendixtitle{А}{Исходный код микросервиса <<Клиенты>>}

% Файл сгенерирован scripts/gen_listings.py — вручную не править

\subsection*{Сборка проекта}

\lstinputlisting[style=xml, style=appendix, caption={pom.xml}]{../../pom.xml}

\subsection*{Общие ресурсы web-клиентов (ui-kit)}

\lstinputlisting[style=xml, style=appendix, caption={ui-kit/pom.xml}]{../../ui-kit/pom.xml}
\lstinputlisting[style=css, style=appendix, caption={ui-kit/src/main/resources/META-INF/resources/ui/bank.css}]{../../ui-kit/src/main/resources/META-INF/resources/ui/bank.css}
\lstinputlisting[style=js, style=appendix, caption={ui-kit/src/main/resources/META-INF/resources/ui/bank.js}]{../../ui-kit/src/main/resources/META-INF/resources/ui/bank.js}

\subsection*{Микросервис «Клиенты» (client-service)}

\lstinputlisting[style=xml, style=appendix, caption={client-service/pom.xml}]{../../client-service/pom.xml}
\lstinputlisting[style=yaml, style=appendix, caption={client-service/src/main/resources/application.yml}]{../../client-service/src/main/resources/application.yml}
\lstinputlisting[style=sql, style=appendix, caption={client-service/src/main/resources/schema.sql}]{../../client-service/src/main/resources/schema.sql}
\lstinputlisting[style=sql, style=appendix, caption={client-service/src/main/resources/data.sql}]{../../client-service/src/main/resources/data.sql}
\lstinputlisting[style=java, style=appendix, caption={client-service/src/main/java/by/bsuir/bank/client/domain/Citizenship.java}]{../../client-service/src/main/java/by/bsuir/bank/client/domain/Citizenship.java}
\lstinputlisting[style=java, style=appendix, caption={client-service/src/main/java/by/bsuir/bank/client/domain/City.java}]{../../client-service/src/main/java/by/bsuir/bank/client/domain/City.java}
\lstinputlisting[style=java, style=appendix, caption={client-service/src/main/java/by/bsuir/bank/client/domain/Client.java}]{../../client-service/src/main/java/by/bsuir/bank/client/domain/Client.java}
\lstinputlisting[style=java, style=appendix, caption={client-service/src/main/java/by/bsuir/bank/client/domain/DictionaryItem.java}]{../../client-service/src/main/java/by/bsuir/bank/client/domain/DictionaryItem.java}
\lstinputlisting[style=java, style=appendix, caption={client-service/src/main/java/by/bsuir/bank/client/domain/Disability.java}]{../../client-service/src/main/java/by/bsuir/bank/client/domain/Disability.java}
\lstinputlisting[style=java, style=appendix, caption={client-service/src/main/java/by/bsuir/bank/client/domain/MaritalStatus.java}]{../../client-service/src/main/java/by/bsuir/bank/client/domain/MaritalStatus.java}
\lstinputlisting[style=java, style=appendix, caption={client-service/src/main/java/by/bsuir/bank/client/domain/Sex.java}]{../../client-service/src/main/java/by/bsuir/bank/client/domain/Sex.java}
\lstinputlisting[style=java, style=appendix, caption={client-service/src/main/java/by/bsuir/bank/client/repository/CitizenshipRepository.java}]{../../client-service/src/main/java/by/bsuir/bank/client/repository/CitizenshipRepository.java}
\lstinputlisting[style=java, style=appendix, caption={client-service/src/main/java/by/bsuir/bank/client/repository/CityRepository.java}]{../../client-service/src/main/java/by/bsuir/bank/client/repository/CityRepository.java}
\lstinputlisting[style=java, style=appendix, caption={client-service/src/main/java/by/bsuir/bank/client/repository/ClientRepository.java}]{../../client-service/src/main/java/by/bsuir/bank/client/repository/ClientRepository.java}
\lstinputlisting[style=java, style=appendix, caption={client-service/src/main/java/by/bsuir/bank/client/repository/DisabilityRepository.java}]{../../client-service/src/main/java/by/bsuir/bank/client/repository/DisabilityRepository.java}
\lstinputlisting[style=java, style=appendix, caption={client-service/src/main/java/by/bsuir/bank/client/repository/MaritalStatusRepository.java}]{../../client-service/src/main/java/by/bsuir/bank/client/repository/MaritalStatusRepository.java}
\lstinputlisting[style=java, style=appendix, caption={client-service/src/main/java/by/bsuir/bank/client/dto/ApiError.java}]{../../client-service/src/main/java/by/bsuir/bank/client/dto/ApiError.java}
\lstinputlisting[style=java, style=appendix, caption={client-service/src/main/java/by/bsuir/bank/client/dto/ClientRequest.java}]{../../client-service/src/main/java/by/bsuir/bank/client/dto/ClientRequest.java}
\lstinputlisting[style=java, style=appendix, caption={client-service/src/main/java/by/bsuir/bank/client/dto/ClientResponse.java}]{../../client-service/src/main/java/by/bsuir/bank/client/dto/ClientResponse.java}
\lstinputlisting[style=java, style=appendix, caption={client-service/src/main/java/by/bsuir/bank/client/dto/Dictionaries.java}]{../../client-service/src/main/java/by/bsuir/bank/client/dto/Dictionaries.java}
\lstinputlisting[style=java, style=appendix, caption={client-service/src/main/java/by/bsuir/bank/client/validation/DateText.java}]{../../client-service/src/main/java/by/bsuir/bank/client/validation/DateText.java}
\lstinputlisting[style=java, style=appendix, caption={client-service/src/main/java/by/bsuir/bank/client/validation/DateTextValidator.java}]{../../client-service/src/main/java/by/bsuir/bank/client/validation/DateTextValidator.java}
\lstinputlisting[style=java, style=appendix, caption={client-service/src/main/java/by/bsuir/bank/client/validation/Dates.java}]{../../client-service/src/main/java/by/bsuir/bank/client/validation/Dates.java}
\lstinputlisting[style=java, style=appendix, caption={client-service/src/main/java/by/bsuir/bank/client/validation/Masks.java}]{../../client-service/src/main/java/by/bsuir/bank/client/validation/Masks.java}
\lstinputlisting[style=java, style=appendix, caption={client-service/src/main/java/by/bsuir/bank/client/ClientServiceApplication.java}]{../../client-service/src/main/java/by/bsuir/bank/client/ClientServiceApplication.java}
\lstinputlisting[style=java, style=appendix, caption={client-service/src/main/java/by/bsuir/bank/client/service/ClientService.java}]{../../client-service/src/main/java/by/bsuir/bank/client/service/ClientService.java}
\lstinputlisting[style=java, style=appendix, caption={client-service/src/main/java/by/bsuir/bank/client/service/FieldErrorException.java}]{../../client-service/src/main/java/by/bsuir/bank/client/service/FieldErrorException.java}
\lstinputlisting[style=java, style=appendix, caption={client-service/src/main/java/by/bsuir/bank/client/service/NotFoundException.java}]{../../client-service/src/main/java/by/bsuir/bank/client/service/NotFoundException.java}
\lstinputlisting[style=java, style=appendix, caption={client-service/src/main/java/by/bsuir/bank/client/web/ApiExceptionHandler.java}]{../../client-service/src/main/java/by/bsuir/bank/client/web/ApiExceptionHandler.java}
\lstinputlisting[style=java, style=appendix, caption={client-service/src/main/java/by/bsuir/bank/client/web/ClientController.java}]{../../client-service/src/main/java/by/bsuir/bank/client/web/ClientController.java}
\lstinputlisting[style=java, style=appendix, caption={client-service/src/main/java/by/bsuir/bank/client/web/DictionaryController.java}]{../../client-service/src/main/java/by/bsuir/bank/client/web/DictionaryController.java}
\lstinputlisting[style=xml, style=appendix, caption={client-service/src/main/resources/static/index.html}]{../../client-service/src/main/resources/static/index.html}
\lstinputlisting[style=js, style=appendix, caption={client-service/src/main/resources/static/app.js}]{../../client-service/src/main/resources/static/app.js}
\lstinputlisting[style=java, style=appendix, caption={client-service/src/test/java/by/bsuir/bank/client/ClientApiIT.java}]{../../client-service/src/test/java/by/bsuir/bank/client/ClientApiIT.java}
\lstinputlisting[style=java, style=appendix, caption={client-service/src/test/java/by/bsuir/bank/client/ClientRequestValidationTest.java}]{../../client-service/src/test/java/by/bsuir/bank/client/ClientRequestValidationTest.java}
\lstinputlisting[style=java, style=appendix, caption={client-service/src/test/java/by/bsuir/bank/client/ClientServiceTest.java}]{../../client-service/src/test/java/by/bsuir/bank/client/ClientServiceTest.java}
\lstinputlisting[style=java, style=appendix, caption={client-service/src/test/java/by/bsuir/bank/client/ClientUiIT.java}]{../../client-service/src/test/java/by/bsuir/bank/client/ClientUiIT.java}
\lstinputlisting[style=java, style=appendix, caption={client-service/src/test/java/by/bsuir/bank/client/DateTextValidatorTest.java}]{../../client-service/src/test/java/by/bsuir/bank/client/DateTextValidatorTest.java}
\lstinputlisting[style=java, style=appendix, caption={client-service/src/test/java/by/bsuir/bank/client/TestClients.java}]{../../client-service/src/test/java/by/bsuir/bank/client/TestClients.java}


\end{document}
